En este reporte se presenta un estudio comparativo de diferentes paradigmas algorítmicos aplicados al problema de optimización combinatoria ``AniMaraton''. Se evaluaron cuatro enfoques: exploración exhaustiva por Fuerza Bruta, dos heurísticas de diseño ávido (Greedy) y una solución óptima mediante Programación Dinámica bidimensional. Los resultados confirman los postulados teóricos: el enfoque exhaustivo sufre un colapso operativo por explosión combinatoria cerca de $N=100$, mientras que la programación dinámica demuestra una escalabilidad pseudo-polinomial eficiente. Las heurísticas Greedy exhibieron tiempos fraccionarios, sacrificando optimalidad en favor de la viabilidad temporal.